\documentclass[11pt, a4paper]{article}

\usepackage{lecture}

\usepackage{notes}
\usepackage{code}
\usepackage{quotes}

\usepackage{shortcuts}

\begin{document}

\makefront{}

\tableofcontents
\newpage

\section{Info}\label{sec:info} % (fold)
\begin{itemize}
    \item \bf{Mario Bravetti} --- \it{mario.bravetti@unibo.it}
    \item \bf{Lorenzo Bacchiani} --- \it{lorenzo.bacchiani2@unibo.it}
\end{itemize}

\paragraph{Esame}\label{par:esame} % (fold)
\begin{itemize}
    \item \bf{Scritto}: $27$ punti
        \begin{warning}
            Il voto dello scritto vale solo nella sessione corrente
        \end{warning}
    \item \bf{Orale di progetto}: $7$ punti (facoltativo)
        Realizzazione di un \bf{compilatore} per un semplice linguaggio di programmazione funzionale \it{object-oriented}, chiamato \it{FOOL}.
        Estensione di quello del \bf{laboratorio}.
        \emptyln{}
        Progetto in singolo o in \bf{gruppo max $4$ persone} e che può essere effettuato \bf{parzialmente} (anche solo un estensione minimale).
        \emptyln{}
        Discussione del progetto con tutto il gruppo su \bf{appuntamento} (dopo lo scritto, entro l'inizio della sessione successiva).
\end{itemize}
% paragraph Esame (end)

% section Info (end)

\section{Linguaggi Regolari}\label{sec:linguaggi_regolari} % (fold)

\subsection{Concetti di Base}\label{sub:concetti_di_base} % (fold)
\begin{definition}{Alfabeto}{alf}
    Insieme finito e non vuoto di simboli
    \begin{itemize}
        \item $\Si = \{0, 1\}$ alfabeto binario
        \item $\Si = \{a, b, c, \dots, z\}$ insieme di tutte le lettere minuscole
        \item Insieme di tutti i caratteri ASCII
    \end{itemize}
\end{definition}

\begin{definition}{Stringa}{str}
    Una sequenza finita di simboli da un alfabeto $\Si$, è una stringa su un alfabeto $\{0, 1\}$.
\end{definition}

\begin{definition}{Stringa vuota}{str_vuota}
    La stringa con zero occorrenze di simboli da $\Si$.
    \begin{itemize}
        \item si denota con $\eps$
    \end{itemize}
\end{definition}

\begin{definition}{Lunghezza di una stringa}{lunghezza_str}
    Numero di posizioni per i simboli nella stringa.
    \begin{itemize}
        \item $\ab\om$ denota la lunghezza della stringa $\om$
    \end{itemize}
\end{definition}

\begin{definition}{Potenze di un alfabeto}{potenze_alf}
    $\Si k \ =$ inisieme delle stringhe di lunghezza $k$ con simboli da $\Si$
    \begin{example}{}{}
        $\Si = \{0, 1\}$ \\
        $\Si[1] = \{0, 1\}$ \\
        $\Si[2] = \{00, 01, 10, 11\}$ \\
        $\Si[0] = \{\eps\}$ \\
        \bf{Domanda}; Quante stringhe ci sono in $\Si[3]$?
        L'insieme di tutte le stringhe su $\Si$ è denotato da $\Si[*]$:
        \[
            \Si[*] = \Si[0] \cup \Si[1] \cup \Si[2] \cup \cdots
        \]
        anche:
        \[
            \begin{matrix}
                \Si[+] = \Si[1] \cup \Si[2] \cup \Si[3] \cup \cdots \\
                \Si[*] = \Si[+] \cup \{\eps\}
            \end{matrix}
        \]
    \end{example}
\end{definition}

\begin{definition}{Concatenazione}{concat_str}
    Se $x$ e $y$ sono stringhe, allora $xy$ è la stringa ottenuta collocando una copia di $y$ subito dopo una copia di $x$.
    \begin{equation*}
        \begin{matrix}
            x = a_1a_2\dots a_i, \quad y = b_1b_2\dots b_j \\
            xy = a_1a_2\dots a_i b_1b_2\dots b_j
        \end{matrix}
    \end{equation*}
    \begin{note}{}{}
        Per ogni stringa $x$:
        \[ x\eps = \eps x = x \]
    \end{note}
\end{definition}

\begin{definition}{Linguaggi}{lang}
    Se $\Si$ è un alfabeto, e $L \in \Si[*]$ allora $L$ è un \bf{linguaggio}.
    \begin{itemize}
        \item $\emptyset \not= \{\eps\}$
        \item L'alfabeto $\Si$ è sempre finito
    \end{itemize}
    \begin{example}{}{}
        \begin{itemize}
            \item L'insieme dele parole italiane legali
            \item L'insieme dei programmi C legali
            \item L'insieme delle stringhe che consistono di $n$ zeri seguiti da $n$ uni:
                \[
                    \{\eps, 01, 0011, 000111, \dots \}
                \]
            \item L'insieme delle stringhe con un numero ugiale di zeri e di uni:
                \[
                    \{\eps, 01, 10, 0011, 0101, 1001, \dots \}
                \]
            \item $L_P =$ insieme dei numeri binari il cui valore è primo:
                \[
                    \{10, 11, 101, 111, 1011, \dots \}
                \]
            \item Il linguaggio vuoto $\emptyset$
            \item Il linguaggio $\{\eps\}$ consiste della stringa vuota
        \end{itemize}
    \end{example}
\end{definition}

% subsection Concetti di Base (end)

\subsection{Automi a stati finiti deterministici}\label{sub:automi_a_stati_finiti_deterministici} % (fold)

\begin{definition}{Automa a stati Finiti Deterministici}{dfa}
Un \bf{DFA} è una quintupla:
    \begin{equation}
        A = \pt{Q, \Si, \dt, q_0, F}
    \end{equation}
    \begin{itemize}
        \item $Q$ è un insieme finito di stati
        \item $\Si$ è un alfabeto finito (simboli in input)
        \item $\dt$ è una funzione di transizione da $Q \times \Si$ a $Q$, cioè: $\pt{q, a}\to p$
        \item $q_0 \in Q$ è lo stato iniziale
        \item $F \subseteq Q$ è un insieme di stati finali
    \end{itemize}
\end{definition}

\begin{example}{Un automa $A$ che accetta $L = \pg{x01y: x, y \in \pg{0, 1}^*}$}{}
    L'automa $\pt{\pg{\qlo, \qli, \qlii}, \pg{0, 1}, \dt, \qlo, \pg \qli}$ definito tramite una \it{tabella di transizione}:
%%% TABELLA
    Lo stesso automa $A$ definito, più semplicemente, tramite un \it{diagramma di transizione}:
%%% IMMAGINE
\end{example}

Un automa a stati finiti (\bf{FA}) \it{accetta} una stringa $w = \ali\alii\dots\aln$ se esiste un cammino nel diagramma di transizione che:
\begin{enumerate}
    \item Inizia nello stato iniziale
    \item Finisce in uno stato finale (di accettazione)
    \item Ha una sequenza di etichette $\ali\alii\dots\aln$
\end{enumerate}

% subsection Automi a stati finiti deterministici (end)

% section Linguaggi Regolari (end)

\end{document}
